% ============================================================================
%  前言
% ============================================================================
\chapter*{前言}
\addcontentsline{toc}{chapter}{前言}
\markboth{前言}{前言}
\thispagestyle{frontstyle}

本书是一份面向量化金融从业者、金融工程专业学生与转行者的系统化学习笔记，覆盖从金融市场基础到随机分析、从衍生品定价到因子投资的完整知识链条。全书 12 章，按「基础—工具—定价—统计—应用」的顺序编排，并在附录中给出公式速查、Python 工具链与术语中英对照。

\section*{为什么需要一份「量化金融工程」的整合笔记}

市面上量化金融的教材与研究文献极其丰富：Hull 的《期权、期货及其他衍生产品》讲清了产品与市场，Shreve 的《金融随机分析》讲透了测度与鞅，Campbell--Lo--MacKinlay 建立了金融计量的方法论，Grinold--Kahn 则系统阐述了主动管理。但在实际学习与工作中，真正的困难往往不是「某一本书读不懂」，而是：

\begin{itemize}[leftmargin=1.8em]
  \item \textbf{知识被割裂}：定价（Q Quant）与统计（P Quant）常被当作两个互不相干的领域，而它们共享同一套概率语言与数值工具；
  \item \textbf{「会用」与「懂」之间的断层}：能调用 Black--Scholes 公式，却说不清为什么风险中性测度存在、为什么波动率不变；
  \item \textbf{理论到工程的距离}：知道 GARCH 的定义，却没有一套可直接运行的代码与检验流程；
  \item \textbf{符号不统一}：不同教材对测度、期望、波动率、久期的记号各不相同，交叉阅读时成本极高。
\end{itemize}

本书的处理方式是：\textbf{用统一的符号体系与统一的叙事逻辑}，把上述内容组织成一条主线——「金融的本质是跨期与跨状态的风险定价；Q Quant 给出无套利约束下的相对价格，P Quant 给出真实测度下的统计规律，而二者都由随机过程与概率统计提供数学基础。」

\section*{内容组织：Q Quant 与 P Quant 的分野与统一}

\begin{description}[leftmargin=2.4em, style=nextline]
  \item[Q Quant（第 6、10 章为核心）] 以风险中性测度 $\Qm$ 为语言，研究「在无套利条件下衍生品应当值多少」。工具是鞅论、测度变换、随机微分方程与数值方法。对应岗位：衍生品定价、结构化产品、模型验证（Model Validation）、XVA。
  \item[P Quant（第 7--9、11 章为核心）] 以真实测度 $\Prob$ 为语言，研究「未来实际会发生什么」，寻找可被统计验证的收益规律。工具是因子模型、时间序列、机器学习与组合优化。对应岗位：因子研究、量化策略、风险管理、组合管理。
\end{description}

需要强调的是：\textbf{二者的区别不是「对的」与「错的」，而是「用途」不同}。给奇异期权报价必须用 $\Qm$（否则可以被无风险套利），做因子择时必须在 $\Prob$ 下检验（因为 $\Qm$ 下的期望收益恒为无风险利率，不含任何 Alpha 信息）。混淆测度是量化面试中的经典陷阱，也是本书反复强调的红线。

\section*{第二版修订说明}

本次修订在保留原有 12 章框架的基础上，做了三方面的系统性升级：

\begin{enumerate}[leftmargin=1.8em]
  \item \textbf{排版重构}：重写统一导言区（\texttt{preamble.tex}），设计封面、统一章节标题样式、页眉页脚与配色方案；补齐目录层级、交叉引用（\texttt{cleveref}）、图表编号与算法环境；修复了旧版本中缺失字形（如警告图标显示为空白）与 PDF 书签乱码的问题。
  \item \textbf{内容扩充}：为每一章补充了公式推导、数值算例、实务应用与风险警示；新增随机微分方程的推导细节、SABR 与局部波动率、XVA 框架、卡尔曼滤波与状态空间模型、已实现波动率与 HAR 模型、Kelly 准则与风险预算、机器学习选股等内容；并为资格证章节补充了可计算的例题。
  \item \textbf{工程化与可视化}：新增 TikZ/pgfplots 绘制的定价与风险示意图（收益曲线、有效前沿、波动率微笑、二叉树、AMM 曲线、布朗路径等），新增可直接运行的 Python 代码片段，新增数学公式附录、Python 工具链附录与术语中英对照表，以及统一整理的参考文献。
\end{enumerate}

\section*{阅读建议}

\begin{itemize}[leftmargin=1.8em]
  \item \textbf{零基础读者}：按第 1--5 章顺序阅读，先建立金融直觉与数学工具，再进入 Q/P Quant 的分野。
  \item \textbf{有金融背景、数学薄弱}：先读第 10 章（随机过程）与第 8 章（概率统计），再回到第 6 章（Q Quant）。
  \item \textbf{求职面试导向}：重点阅读第 4、6、7、11 章的公式与「实务应用」盒子，第 12 章了解美国证券从业资格体系；面试高频问题集中在测度变换、伊藤引理、希腊字母、因子模型与回测陷阱。
  \item \textbf{工程导向}：直接看附录 B（Python 工具链）与各章的代码示例，再回溯对应的数学推导。
\end{itemize}

\section*{关于符号、代码与免责声明}

\begin{itemize}[leftmargin=1.8em]
  \item 全书数学符号尽量统一，完整列表见「符号表」；不同文献的记号差异在相应位置会给出对照。
  \item 代码示例以 Python 生态为主（\texttt{numpy}、\texttt{pandas}、\texttt{statsmodels}、\texttt{arch}、\texttt{scipy}、\texttt{cvxpy} 等），为便于跨平台运行，变量名与注释一律使用英文。
  \item 本书是学习笔记的整理，\textbf{不构成任何投资建议}。所有模型的假设与局限都已在相应章节的「风险警示」中尽量说明；历史回测结果不代表未来收益。
\end{itemize}

\vspace{1.2em}
\begin{flushright}
  \sffamily Quantex\\
  \sffamily \today
\end{flushright}

\clearpage
